\chapter{Diseño e implementación del modelo a desarrollar}
\label{cap:diseno}

Este capítulo articula el modelo de investigación y su materialización en software. Se presenta primero el \emph{diseño metodológico} —el tipo de investigación y el modelo de simulación numérica que sostiene la verificación, las variables y su matriz de consistencia, la población, los casos de estudio y su criterio de selección, el procedimiento con sus instrumentos numéricos y documentales, y los criterios de aceptación con que se contrasta la hipótesis— y, a continuación, el software EduFEM, que concreta ese modelo.

\section{Diseño metodológico}
\label{sec:metodologia}

La Introducción presentó una visión general del enfoque de investigación; esta sección lo desarrolla siguiendo la estructura de un modelo de investigación: la conceptualización y el alcance del modelo, la definición de las variables, el procedimiento y la toma de observaciones, la validación de esas observaciones y el contraste de la hipótesis con los resultados.

\subsection{Tipo de investigación y modelo de simulación numérica}
\label{sec:proceso-met}

Por su finalidad, el trabajo es de carácter \emph{propositivo} —la propuesta de un artefacto que resuelve un problema formativo—; por su naturaleza es investigación aplicada de tipo \emph{tecnológico}, la que produce conocimiento operativo para transformar una realidad concreta \autocite[pp.~90-92]{garciacordoba2005tecnologica}; por el abordaje del problema, el enfoque es \emph{cuantitativo} en la verificación y validación del motor \autocite[p.~6]{sampieri2018metodologia}, pues la corrección del artefacto se establece midiendo errores numéricos, tensiones, desplazamientos y tasas de convergencia frente a referencias de exactitud conocida, y \emph{documental} en la validación de contenido, que confronta cada instrumento del software con las destrezas que un consenso publicado de expertos exige y cada decisión de diseño con un principio de aprendizaje; y por su nivel es \emph{descriptiva} y \emph{comparativa}: caracteriza el comportamiento numérico del motor, contrasta de forma sistemática el desempeño de los elementos Q4 y Q9 y describe el grado en que el artefacto expone lo que la hipótesis postula. El modelo que sostiene la verificación es un \emph{modelo de simulación numérica} —una representación ideal del objeto en la que el investigador abstrae y sistematiza las relaciones que considera esenciales \autocite[p.~21]{alvarez_metodologia}, ejercitada aquí por simulación computacional y no por experimentación física—: el análisis de medios continuos en elasticidad lineal plana por el MEF, que es el contenido del objeto de estudio, se ejercita resolviendo problemas de ese dominio con el motor desarrollado y contrastando sus resultados con referencias de exactitud conocida, de modo que la respuesta que el software expone al estudiante sea correcta. No hay unidades experimentales, tratamientos ni aleatorización: el motor es determinista, de modo que la confiabilidad del modelo no descansa en réplicas sino en la reproducibilidad exacta de cada corrida y en la convergencia de sus resultados bajo refinamiento de malla; la incertidumbre numérica se cuantifica con las tasas de convergencia y, cuando la referencia no es exacta, con la extrapolación de Richardson de la propia secuencia de mallas \autocite[pp.~309-312]{oberkampf2010vv}. La unidad de análisis de esta etapa es el artefacto: las variables del problema científico se observan en él con los instrumentos de la \autoref{sec:validacion-datos}, y su efecto sobre el estudiante se deja a la prueba de campo que las recomendaciones diseñan.

El trabajo combinó dos ciclos articulados. El \emph{desarrollo del software} siguió un proceso iterativo e incremental: cada componente del canal de cálculo del MEF (\autoref{cap:marco_teorico}) se implementó y se sometió a prueba de regresión antes de integrarse al conjunto. La \emph{evaluación} se realizó en dos frentes: la corrección del motor, mediante verificación de código por el método de soluciones manufacturadas y contrastación con dos referencias externas, la solución analítica de la viga de Timoshenko y la membrana de Cook \autocite{oberkampf2010vv,salari2000mms}; y la validez de contenido y coherencia del artefacto, mediante la tabla de especificaciones y la matriz de principios de la \autoref{sec:validacion-datos}, construidas sobre la versión del software que acompaña a este documento.

\subsection{Variables e indicadores}
\label{sec:variables-met}

La \autoref{tab:variables} operacionaliza las variables del estudio, es decir, especifica para cada una las operaciones concretas con que se la mide \autocite[p.~137]{sampieri2018metodologia}. La \emph{variable independiente} es la forma en que el software expone el canal de cálculo y la respuesta estructural: como caja negra —solo datos de entrada y resultados— o de manera transparente, interactiva y numéricamente verificable. Se observa en el artefacto con indicadores contables: etapas del canal con módulo interactivo y con desarrollo numérico en la memoria, resultados expuestos en el post-proceso con su procedimiento, fases que comparten el lienzo, conmutador entre fórmula y valor, y coincidencia entre lo que el software expone y lo que su motor verificado calcula. La \emph{variable dependiente} es el criterio del estudiante para interpretar y juzgar la respuesta estructural obtenida por el MEF: desplazamientos, tensiones, reacciones, calidad de la malla y fenómenos numéricos. De ella, esta etapa observa la dimensión que puede establecerse sobre el artefacto —su \emph{validez de contenido}, medida por la proporción de destrezas de interpretación y verificación exigidas por los expertos que tienen instrumento en el software, y su \emph{coherencia}, medida por la proporción de principios de aprendizaje con encarnación comprobable— y deja a la prueba de campo la dimensión que reside en el estudiante, el \emph{criterio logrado}, que se mediría con la ganancia entre una preprueba y una posprueba de juicio de resultados. La \emph{variable interviniente} es EduFEM, la solución tentativa con que se manipula la independiente. Las magnitudes del experimento numérico con que se establece la corrección del motor no son variables del problema científico, y la tabla las distingue con ese nombre: los \emph{factores} —la definición del modelo estructural (geometría, material, acciones y apoyos) y las decisiones de discretización, el tipo de elemento (Q4 o Q9) y la densidad de malla (tamaño característico $h$, o $N\times N$)— y las \emph{magnitudes medidas} —la respuesta estructural calculada (tensiones y desplazamientos) junto con su exactitud frente a las soluciones de referencia, medida por las normas de error $L^2$ y $H^1$ del desplazamiento, por la norma $L^2$ del campo de tensiones recuperado y por los errores relativos en tensión normal y en flecha—. El tipo de análisis (tensión o deformación plana) y el orden de cuadratura se mantienen como \emph{controlados}, y el desempeño del solucionador se registra como \emph{observado}, sin criterio de aceptación, para documentar la escalabilidad del motor.

{\footnotesize
\setlength{\tabcolsep}{3pt}
\renewcommand{\arraystretch}{1.15}
\begin{longtable}{@{}>{\raggedright\arraybackslash}p{2.9cm}>{\raggedright\arraybackslash}p{3.0cm}>{\raggedright\arraybackslash}p{5.0cm}>{\raggedright\arraybackslash}p{2.7cm}@{}}
\caption{Operacionalización de las variables de la investigación.}\label{tab:variables}\\
\toprule
\textbf{Variable} & \textbf{Dimensión} & \textbf{Indicador} & \textbf{Instrumento} \\
\midrule
\endfirsthead
\caption[]{Operacionalización de las variables de la investigación (continuación).}\\
\toprule
\textbf{Variable} & \textbf{Dimensión} & \textbf{Indicador} & \textbf{Instrumento} \\
\midrule
\endhead
\bottomrule
\endfoot
Independiente: exposición del canal de cálculo y de la respuesta & Transparencia; interactividad; verificabilidad numérica & Etapas con módulo interactivo y con desarrollo numérico en la memoria; resultados expuestos con su procedimiento en el post-proceso; fases que comparten el lienzo; conmutador fórmula--valor; coincidencia entre lo expuesto y lo calculado por el motor verificado & Inventario de módulos; memoria del ejemplo canónico; pruebas de generación de la memoria \\
Dependiente: criterio para interpretar la respuesta estructural & Validez de contenido & Destrezas de interpretación y verificación exigidas por los expertos con instrumento en el software & Tabla de especificaciones \\
 & Coherencia del diseño & Principios de aprendizaje con encarnación comprobable & Matriz de principios \\
 & Criterio logrado (no medido en esta etapa) & Ganancia entre preprueba y posprueba de juicio de resultados & Prueba de campo (recomendaciones) \\
Interviniente: EduFEM & Versión evaluada & Software distribuido con este documento & Instalador y repositorio \\
\midrule
Factores del experimento numérico & Definición del modelo; discretización & Dominio y malla; $E$, $\nu$, espesor $t$; cargas y restricciones; Q4/Q9; tamaño característico $h$ (o $N\times N$) y número de GDL & Guion de cada caso \\
Magnitudes medidas & Respuesta estructural; exactitud & $\sigma_x,\sigma_y,\tau_{xy},\sigma_{\mathrm{VM}}$; $u,v$ nodales; tasa de convergencia (orden $p$); normas $L^2$ y $H^1$ del desplazamiento y $L^2$ de $\sigma^*$; error relativo (\%) en $\sigma_x$ y en flecha frente al analítico y a SAP2000; residuo de la suma de reacciones & Solucionador del MEF; guiones de V\&V \\
Controlados & Tipo de análisis; integración & Tensión o deformación plana; orden de cuadratura & Constantes del guion de cada caso \\
Observado: desempeño & Tiempo; memoria & Tiempos de ensamblaje y de solución; memoria de $\bm{K}$ densa y dispersa & Guion de medición de tiempos \\
\end{longtable}}

\subsection{Matriz de consistencia}
\label{sec:matriz-consistencia}

El \textbf{problema científico}, en la formulación única de la Introducción —¿de qué manera el uso del software de análisis estructural como caja negra podrá afectar en el bajo criterio para interpretar la respuesta estructural obtenida por el MEF en los estudiantes de la Carrera de Ingeniería Civil de la Universidad Autónoma Tomás Frías en la gestión 2026?—, se atiende mediante el \textbf{objetivo general} de desarrollar EduFEM y se contrasta con la \textbf{hipótesis} de que la exposición transparente e interactiva del canal y de la respuesta permitirá mejorar ese criterio, cuya parte comprobable en esta etapa son las tres cláusulas de exposición, corrección y validez de contenido. La \autoref{tab:consistencia} traza la correspondencia entre cada objetivo específico, la pregunta de investigación que responde, las variables o indicadores asociados, el método o instrumento empleado y la evidencia que lo respalda en el documento.

\begin{table}[htbp]
\centering
\caption[Matriz de consistencia de la investigación.]{Matriz de consistencia: correspondencia entre objetivos específicos, preguntas de investigación, variables, métodos y evidencia.}
\label{tab:consistencia}
{\footnotesize
\setlength{\tabcolsep}{3pt}
\renewcommand{\arraystretch}{1.15}
\begin{tabular}{@{}>{\raggedright\arraybackslash}p{3.0cm}>{\raggedright\arraybackslash}p{1.6cm}>{\raggedright\arraybackslash}p{3.1cm}>{\raggedright\arraybackslash}p{3.3cm}>{\raggedright\arraybackslash}p{2.6cm}@{}}
\toprule
\textbf{Objetivo específico} & \textbf{Pregunta asociada} & \textbf{Variable / indicador} & \textbf{Método / instrumento} & \textbf{Evidencia} \\
\midrule
Diagnosticar la brecha entre operar el software y juzgar su respuesta & PI-1 a PI-3 & Atributos que la herramienta debe tener & Revisión de los antecedentes y del estado del arte; análisis comparativo & \autoref{sec:antecedentes}; \autoref{tab:comparativa} \\
Fundamentar el MEF en elasticidad plana y los principios de aprendizaje & PI-1 a PI-3 & Canal de cálculo cubierto; principios adoptados & Revisión de la literatura clásica del MEF, de la V\&V y de la enseñanza del método & \autoref{cap:marco_teorico} \\
Desarrollar el motor, la interfaz, los módulos y la memoria & PI-1 & Etapas con módulo y con desarrollo en la memoria; resultados expuestos; fases que comparten el lienzo; formatos con ida y vuelta & Implementación en NumPy/SciPy; capas interactivas sobre la malla real; generación de PDF con las funciones del motor & \S\ref{sec:motor}--\S\ref{sec:memoria-io}; \autoref{tab:fases}; \autoref{sec:consistencia-interna} \\
Verificar y validar el motor & PI-2 & Exactitud; tasa de convergencia; bloqueo por cortante & MMS, viga de Timoshenko, membrana de Cook & \autoref{cap:resultados} \\
Validar el contenido y la coherencia del software & PI-3 & Destrezas con instrumento; principios con encarnación & Tabla de especificaciones; matriz de principios & \autoref{tab:especificaciones}; \autoref{tab:principios}; \autoref{sec:validacion-hipotesis} \\
\bottomrule
\end{tabular}}
\end{table}

Las preguntas de investigación referidas en la matriz son las enunciadas en la Introducción, una por cláusula de la hipótesis: \textbf{PI-1}, sobre la organización de la interfaz, de los módulos y de la memoria que expone el canal y la respuesta en lugar de encapsularlos; \textbf{PI-2}, sobre la corrección del motor frente a soluciones analíticas y comerciales y la reproducción de los fenómenos numéricos; y \textbf{PI-3}, sobre la validez de contenido frente a las destrezas que los expertos exigen y la coherencia con los principios de aprendizaje. Los dos primeros objetivos específicos —el diagnóstico y la fundamentación— son transversales: sustentan las tres preguntas y se materializan en el \autoref{cap:marco_teorico}, de ahí que la matriz les asocie las tres. La interoperabilidad de datos (DXF, CSV), incluida en el tercer objetivo, es complemento instrumental: no deriva de una pregunta de investigación, pero se verifica por ida y vuelta e idempotencia (\autoref{sec:consistencia-interna}) para que el modelo del alumno pueda entrar y salir de la herramienta sin pérdida.

\subsection{Población, universo de contenido y casos de estudio}
\label{sec:poblacion}

La \emph{población} del problema científico son los estudiantes de la Carrera de Ingeniería Civil de la Universidad Autónoma Tomás Frías que cursan las asignaturas en que se introduce el análisis por el MEF. % DATO PENDIENTE: nombre, sigla y semestre de la asignatura segun el plan de estudios vigente
Esta etapa no toma una muestra de esa población, porque su unidad de análisis es el artefacto: el criterio del estudiante no se mide, y lo que se observa en él queda diferido a la prueba de campo de las recomendaciones, cuya muestra será dirigida —un curso completo, seleccionado por accesibilidad y no por procedimiento probabilístico \autocite[p.~215]{sampieri2018metodologia}—. En su lugar, esta etapa define un \emph{universo de contenido} que recorre de forma exhaustiva: las destrezas de interpretación y verificación de la respuesta y los conceptos del canal de cálculo que el consenso de 67 expertos de la industria y de la academia reunido por Pérez-Santiago y Campos califica como necesarios para el egresado \autocite[pp.~1162-1163]{perezsantiago2023fem}. De sus 49 ítems se toman los que definen el criterio sobre la respuesta y que caen dentro del alcance de la elasticidad plana: los seis de mallado, los tres de post-proceso, los tres de conceptos fundamentales y los tres de verificación y validación —quince destrezas, agrupadas en las dimensiones que ese consenso califica más alto \autocite[p.~1168]{perezsantiago2023fem}— más los tres conceptos teóricos del canal que el software expone: el ensamblaje con restricciones y solución, la formulación isoparamétrica de elementos bidimensionales y la integración numérica de Gauss. Los ítems restantes —análisis tridimensional, placas y cáscaras, térmico, CAD, tipos de análisis más allá del estático lineal— quedan fuera del alcance declarado y se listan en la \autoref{sec:contenido-didactico} para que la exclusión sea explícita. A ese universo se suman los cuatro principios de aprendizaje de la \autoref{sec:fundamentos-pedagogicos}, que la matriz de principios recorre uno por uno.

Para el experimento numérico, la población sería el universo de problemas de elasticidad lineal plana —en tensión o deformación plana— resolubles con elementos cuadriláteros isoparamétricos. El conjunto de casos que sigue no es una muestra estadística de ese universo, sino una selección \emph{intencional} por valor probatorio —una muestra dirigida, no probabilística \autocite[p.~215]{sampieri2018metodologia}—, coherente con la práctica de la verificación de códigos de cálculo, donde los casos de prueba se eligen por su capacidad de ejercitar todos los términos del modelo matemático y por cubrir la topología de malla más general que el código deberá resolver, antes que por su realismo físico \autocite{oberkampf2010vv}. Cada caso aporta una referencia de naturaleza distinta y ejercita términos distintos del modelo. El método de soluciones manufacturadas provee una solución exacta arbitraria y ejercita la fuerza de volumen, las restricciones de Dirichlet no homogéneas, las dos matrices constitutivas y el mapeo sobre elementos distorsionados, en cuatro configuraciones (cuadrado unitario y cuadrilátero distorsionado, en tensión y en deformación plana). La viga de Timoshenko simplemente apoyada ejercita la carga superficial uniforme y los apoyos, y aporta las dos referencias externas: una solución analítica de la teoría de la elasticidad y un modelo de cáscara en SAP2000. La membrana de Cook ejercita la distorsión trapezoidal bajo cortante y flexión, el modo de deformación ante el que el Q4 bloquea. Entre los tres quedan cubiertos los términos del modelo de la \autoref{sec:formulacion-debil} —cargas nodales, volumétricas y superficiales uniformes, restricciones homogéneas y no homogéneas, ambos estados planos y ambos elementos—, con una excepción que se declara: la carga superficial linealmente variable no interviene en ningún caso de referencia y no tiene contraste externo; las recomendaciones proponen un caso que la incorpore. El refinamiento de malla y la comparación Q4 frente a Q9 se realizan en los dos casos con secuencia de mallas —el MMS y la membrana de Cook, ambos con $N\in\{2,4,8,16,32\}$—; la viga de Timoshenko se resuelve con una única malla fina de elementos Q9 y constituye un contraste puntual de exactitud frente a las dos referencias externas, no un estudio de convergencia.

\subsection{Procedimiento, instrumentos y reproducibilidad}
\label{sec:validacion-datos}

Cada caso numérico sigue el mismo procedimiento. Un guion construye la malla y define material, cargas y apoyos; un validador de salud del modelo —una función pura sin dependencias de la interfaz— verifica la consistencia de los datos de entrada antes de resolver: existencia de elementos, restricciones suficientes para suprimir los movimientos de cuerpo rígido, ausencia de nodos referenciados inexistentes y no degeneración de los elementos (determinante Jacobiano positivo); los errores críticos bloquean el cálculo. En la verificación por soluciones manufacturadas la consistencia de los datos está garantizada por construcción: el término fuente se deduce analíticamente de la solución exacta elegida a priori, de modo que las cargas y las restricciones de Dirichlet corresponden exactamente a ella. El sistema se ensambla y resuelve con las mismas funciones del motor que consumen los módulos educativos y la memoria de cálculo, lo que garantiza coherencia entre lo que el alumno ve explicado y lo que el programa resuelve, y las tensiones se recuperan por la secuencia descrita en el \autoref{cap:marco_teorico}.

Los instrumentos de medición numérica son los guiones de verificación y validación del proyecto (\texttt{tests/vv\_mms.py}, \texttt{tests/vv\_timoshenko.py} y \texttt{tests/vv\_cook.py}), que computan las normas de error —integradas con un punto de Gauss más por dirección que la rigidez (\autoref{sec:verificacion-validacion}) para no subestimar artificialmente el error—, los errores relativos frente a las referencias y las tasas de convergencia, y depositan los datos crudos en archivos CSV. Dentro de cada caso solo se manipulan el tipo de elemento y la densidad de malla —fijos en la viga de Timoshenko, por lo dicho en la \autoref{sec:poblacion}—; el tipo de análisis, el material, la geometría, las cargas y los apoyos, el orden de cuadratura ($2\times2$ para Q4 y $3\times3$ para Q9) y las tolerancias numéricas del motor —el cero numérico y el determinante Jacobiano mínimo admisible, constantes del programa— permanecen fijos.

Los instrumentos documentales son dos y se construyen sobre la versión del software que acompaña a este documento, con las funciones que la \autoref{sec:diseno-edufem} describe y las capturas que el \autoref{anexo:manual} reproduce. El primero es la \emph{tabla de especificaciones}, la vía documental de la validez de contenido: cruza cada destreza y concepto del universo de la \autoref{sec:poblacion} con el instrumento del software que lo cubre y con la evidencia de este documento que lo demuestra; una celda se marca solo si el instrumento existe en el software distribuido y el \autoref{cap:diseno} lo describe, y queda vacía en caso contrario. El criterio externo es el consenso publicado de expertos, que hace las veces de juicio de expertos sin panel propio: la limitación de ese criterio —expertos mexicanos de ingeniería mecánica, que el propio estudio declara no generalizables a otros contextos \autocite[p.~1171]{perezsantiago2023fem}— se acota porque los ítems que se toman son contenido universal del método y no preferencias de un país. El segundo es la \emph{matriz de principios}: para cada uno de los cuatro principios de aprendizaje de la \autoref{sec:fundamentos-pedagogicos} registra la fuente, la decisión de diseño de la \autoref{sec:diseno-edufem} que lo encarna y la evidencia comprobable en el software —una función descrita en ese capítulo, una captura del manual o una prueba de la batería de regresión—. A estos dos instrumentos se suma el análisis comparativo con los antecedentes de la \autoref{tab:comparativa}, que sitúa el artefacto frente a los recursos existentes y sostiene el primer objetivo específico.

La reproducibilidad se apoya, en lo numérico, en el determinismo del motor: ante las mismas entradas produce las mismas salidas, y cada caso se regenera por completo al reejecutar su guion con las dependencias declaradas en el repositorio. A ello se suma una batería de pruebas de regresión que verifica la estabilidad del modelo de datos ante operaciones de edición —el ciclo de transformación $\text{Q4}\rightarrow\text{Q9}\rightarrow\text{Q4}$ sin deriva numérica, la equivalencia de resultados con identificadores de nodo no contiguos y la interoperabilidad de formatos— y que se reporta en el \autoref{cap:resultados}. En lo documental, la reproducibilidad descansa en que cada celda de las dos tablas remite a una función descrita en la \autoref{sec:diseno-edufem}, a una captura del \autoref{anexo:manual} o a una prueba de la batería de regresión, de modo que cualquier lector puede rehacer cada veredicto con el software distribuido y el manual; su limitación, que se declara, es que las construye el autor.

\subsection{Criterios de aceptación y contraste de la hipótesis}
\label{sec:validacion-resultados}

Los resultados numéricos se contrastan con referencias de naturaleza distinta y con criterios fijados a priori en los propios guiones, que los evalúan automáticamente y terminan con error si alguno falla. Para el motor: tasas de convergencia del desplazamiento observadas dentro de $\pm0{,}5$ de las teóricas ($\mathcal{O}(h^{p+1})$ en $L^2$ y $\mathcal{O}(h^{p})$ en $H^1$) en las cuatro configuraciones del MMS; para el campo de tensiones recuperado, cuya tasa teórica no está establecida para la secuencia de extrapolación y promediado nodal que emplea el motor, una cota empírica mínima declarada como tal —tasa no inferior a $1{,}4$ en Q4 ni a $1{,}9$ en Q9—; en la viga de Timoshenko, error de la flecha central inferior al $3\,\%$ frente a la solución analítica con corrección de cortante, error de la tensión normal $\sigma_x$ inferior al $1\,\%$ frente a la solución analítica y frente a SAP2000 en los tres puntos de control, y equilibrio global —suma de reacciones verticales igual a la carga total— con residuo relativo inferior a $10^{-8}$; y, en la membrana de Cook, desplazamiento del Q9 con $N=8$ dentro del $1{,}5\,\%$ del valor de referencia y flecha del Q4 inferior a la del Q9 con la misma malla, que es la firma del bloqueo por cortante. El desempeño del solucionador se reporta sin criterio de aceptación. El \autoref{cap:resultados} reporta los valores efectivamente medidos en cada magnitud, incluidas las componentes secundarias de tensión que no forman parte del criterio.

Para la exposición: módulo interactivo para cada una de las siete etapas elementales y de ensamblaje del canal (M1 a M7) y desarrollo con sustitución numérica de las nueve etapas en la memoria de cálculo; la solución y la recuperación de tensiones expuestas en el post-proceso con sus modos crudo y suavizado; las tres fases del análisis operando sobre un mismo lienzo; ida y vuelta CSV/ZIP sin pérdida de entidades ni de resultados, reimportación DXF idempotente y generación de la memoria para Q4 y Q9 sin error. Para la validez de contenido y la coherencia: cada una de las quince destrezas y de los tres conceptos del universo tiene al menos un instrumento identificable en el software, y cada uno de los cuatro principios de aprendizaje tiene al menos una encarnación comprobable; una destreza sin instrumento o un principio sin encarnación refutarían la cláusula.

La hipótesis enunciada en la Introducción se contrasta cláusula por cláusula con esos criterios: la cláusula (a), exposición, con los indicadores de cobertura, fases y formatos (PI-1); la cláusula (b), corrección, con los criterios numéricos del motor (PI-2); y la cláusula (c), validez de contenido y coherencia, con la tabla de especificaciones y la matriz de principios (PI-3). El efecto de la exposición sobre el criterio del estudiante no se contrasta en esta etapa: se fundamenta con la evidencia de la \autoref{sec:fundamentos-pedagogicos} y su prueba de campo se diseña en las recomendaciones con la población de la \autoref{sec:poblacion}.
